\begin{enumerate}
	\item Is the ethernet cable plugged in (are laps blinking) / are you connected to wifi?
	\item Do you know the IP address of the robot?
\begin{itemize}
	\item Do you know your own IP address on LAN?
		Use nmcli, ifconfig, or ip addr to find it. 
		Also note the interface for this IP. Is it ethernet or wifi? Do you know which one you want?
	\item Find all IPs on a network with arp-scan. 
		You might have to specify the network interface to get output (ex. eth0). 
		You also might have to run this as root (put sudo in front of command)
		Example:
		\newline
		arp-scan 192.168.0.1/24 -I eth0
\end{itemize}
\item Is the robot actually connected to the network? Send it a ping with ping ROBOT\_IP.
\item If on Docker, are you sure you share your host network with the Docker container?
	Do all network checks in this list on both your host and inside the Docker container.
	If your OS does not support --net=host on Docker, make sure you're forwarding the necessary ports.
\item Are you sure the robot is ready to receive commands?
\begin{itemize}
\item Is the robot set to "remote control" or whatever similar setting for you to be able to connect to it?
\item Is the robot in "emergency stop mode"?
\item Are the robot's motor on/initialized?
\end{itemize}
\item Are you sure can connect to the robot? 
\begin{itemize}
	\item Try ssh-ing into it/ open its dashboard in WEB/whatever else to make sure connection is fine, and not some problem with 
		the control port, or the driver. 
\end{itemize}
\item Are you sure can use the remote-control interface/robot driver? 
\begin{itemize}
	\item Did the previous interface exit cleanly? The socket might not have been cleaned up. Try running again. 
		Look at the terminal for specific errors. Try running again.
	\item If using ROS, are you sure all topics are brought up? Do topic list. Then do topic echo on all of them.
\end{itemize}
\item Remove all bells and whistles to test just connection to the robot, maybe just one of the following is the problem:
\begin{itemize}
	\item Skip connecting to the gripper, if any.
	\item Skip parts of the robot. Say on a mobile manipulator, connect to just the arm. Then just the base. 
	\item Skip initializing f/t sensors, if any.
	\item Skip initializing the camera, if any.
	\item If on ROS, do the smallest possible bringup.
	\item Use the simplest possible control example (I always do a point-to-point clik).
\end{itemize}
\item Reboot the robot, run through all the steps again.
\item Reboot your computer, run through all the steps again.
\item Use a different computer, run through all the steps again.
\item Are you using a 1-1 to connection to the robot, or is it going through some LAN in your lab? 
	This intermediary might be the culprit.
	Switch to a 1-1 connection (network is just your computer and the robot, create new fixed connection),
	i.e. plug the ethernet cable directly from the robot into your computer.
	If the robot itself has a router, switch etc, try skipping over it and go directly to the robot.
	Maybe this router/something else has to run DHCP, and you have to let it assign you an IP.
	Try both a fixed IP and let whatever assign the IP to you. The same goes for the robot itself.
	(had a gripper that had to run it's DHCP, and everything obey the gripper).
\item Disable (all) firewall(s) on your computer, if any, try again 
	(had a Windows computer which did not allow socket.bind() period, due to idk what).
\item Maybe something else is preventing socket binding/listening. ssh into the robot. Try connecting with netcat. 
	nc -kl 0.0.0.0 7777 on robot, try nc ROBOT\_IP 7777 on your computer. This should be an echo server.
	Try nc -kl ROBOT\_IP 7777 on robot also. Try using different ports.
	Make sure it works the other way around, i.e. nc -kl on your machine , nc from robot (-l is listening, meaning that's the server).
\item If you got to here, you can connect to the robot and the interface doesn't work for whatever reason.
	Try putting your control code onto the robot's computer, run it from there. First just try the driver from the robot itself.
\item Send an email/open a github issue to whoever gave you/wrote the interface to the robot, give them all the error logs
	and convince them it's not something in the connection or something preventing connections on your computer.
\end{enumerate}
